\documentclass[CJK]{beamer}
\usepackage{CJKutf8}
\usepackage{beamerthemesplit}
\usetheme{Malmoe}
\useoutertheme[footline=authortitle]{miniframes}
\usepackage{amsmath}
\usepackage{amssymb}
\usepackage{graphicx}
\usepackage{color}
\usepackage{slashed}
\usepackage{simplewick}
\graphicspath{{../figures/}}
\def\be{\begin{equation}}
\def\ee{\nonumber\end{equation}}
\def\bea{\begin{eqnarray}}
\def\eea{\nonumber\end{eqnarray}}
\def\ii{{\dot{\imath}}}
\def\bch{\begin{CJK}{UTF8}{gbsn}}
\def\ech{\end{CJK}}
\def\bex{\begin{minipage}{0.3\textwidth}\includegraphics[width=1in]{jugelizi.png}\end{minipage}\begin{minipage}{0.6\textwidth}}
\def\eex{\end{minipage}}
\def\chtitle#1{\frametitle{\bch#1\ech}}
\def\skipline{{\vskip0.1in}}
\def\skiplines{{\vskip0.2in}}
\def\lagr{{\mathcal{L}}}
\def\hamil{{\mathcal{H}}}
\def\vecv{{\mathbf{v}}}
\def\vecx{{\mathbf{x}}}
\def\veck{{\mathbf{k}}}
\def\vecp{{\mathbf{p}}}
\def\vecn{{\mathbf{n}}}
\def\vecA{{\mathbf{A}}}
\def\vecP{{\mathbf{P}}}
\def\vecsigma{{\mathbf{\sigma}}}
\def\hatJn{{\hat{J_\vecn}}}
\def\hatJx{{\hat{J_x}}}
\def\hatJy{{\hat{J_y}}}
\def\hatJz{{\hat{J_z}}}
\def\hatj#1{\hat{J_{#1}}}
\def\hatphi{{\hat{\phi}}}
\def\hatq{{\hat{q}}}
\def\hatpi{{\hat{\pi}}}
\def\vel{\upsilon}
\def\Dint{{\mathcal{D}}}
\def\adag{{\hat{a}^\dagger}}
\def\bdag{{\hat{b}^\dagger}}
\def\cdag{{\hat{c}^\dagger}}
\def\ddag{{\hat{d}^\dagger}}
\def\hata{{\hat{a}}}
\def\hatb{{\hat{b}}}
\def\hatc{{\hat{c}}}
\def\hatd{{\hat{d}}}
\def\hatN{{\hat{N}}}
\def\hatH{{\hat{H}}}
\def\hatp{{\hat{p}}}
\def\Fup{{F^{\mu\nu}}}
\def\Fdown{{F_{\mu\nu}}}
\def\newl{\nonumber \\}
\def\SIkm{\mathrm{km}}
\def\SIyr{\mathrm{yr}}
\def\SIGyr{\mathrm{Gyr}}
\def\SIeV{\mathrm{eV}}
\def\SIGeV{\mathrm{GeV}}
\def\SIm{\mathrm{m}}
\def\SIcm{\mathrm{cm}}
\def\SIJ{\mathrm{J}}
\def\SIs{\mathrm{s}}
\def\SIkg{\mathrm{kg}}
\def\SIg{\mathrm{g}}
\def\vece{\mathrm{e}}
\def\bmat#1{\left(\begin{array}{#1}}
\def\emat{\end{array}\right)}
\def\bcase#1{\left\{\begin{array}{#1}}
\def\ecase{\end{array}\right.}
\def\calM{{\mathcal{M}}}
\def\calT{{\mathcal{T}}}
\def\calR{{\mathcal{R}}}
\def\barpsi{\bar{\psi}}
\def\baru{\bar{u}}
\def\barv{\bar{\upsilon}}
\def\bmini#1{\begin{minipage}{#1\textwidth}}
\def\emini{\end{minipage}}
\def\qeq{\stackrel{?}{=}}
\def\torder#1{\mathcal{T}\left(#1\right)}
\def\rorder#1{\mathcal{R}\left(#1\right)}


\title{Quantum Field Theory I \\ Homework 6}
  \author{}
  \date{}


\begin{document}

\begin{frame}
 
\begin{center}
\begin{Large}
\bch
量子场论 I 

{\vskip 0.3in}

第七次 课后作业 （5题共2.5分, 模拟期末考试题）

交作业时间: 12月19日，星期一，13：30pm

\ech
\end{Large}
\end{center}

\vskip 0.2in

\bch
课件下载
\ech
https://github.com/zqhuang/SYSU\_QFTI

\end{frame}


\begin{frame}
\chtitle{第1题(0.5分）}
\bch
如果不改变时空的度规，是否可能从真空中产生出一对正负粒子？为什么？
\ech
\end{frame}

\begin{frame}
\chtitle{第2题(0.5分）}
\bch
阐述什么是自由场。为什么自由场量子化之前先要对场进行傅立叶变换？
\ech
\end{frame}

\begin{frame}
\chtitle{第3题(0.5分）}
\bch
在课上我们把自由实标量场量子化为
$$ \hatphi(x) = \frac{1}{(2\pi)^{3/2}} \int \sqrt{\frac{d^3\veck}{2\omega}} \left(\hata_{\veck} e^{-ik_\mu x^\mu} + \adag_{\veck}e^{ik_\mu x^\mu}\right) $$

考虑两个实标量场$\phi$和$\psi$，拉氏密度为
$$\lagr = \frac{1}{2}\partial^\mu\phi\partial_\mu\phi + \frac{1}{2}\partial^\mu\psi\partial_\mu\psi - \frac{1}{2}m^2(\phi^2+\psi^2+\phi\psi) $$
把$\phi$和$\psi$都量子化 (写成独立的产生湮灭算符的线性迭加)。
\ech
\end{frame}


\begin{frame}
\chtitle{第4题(0.5分）}
\bch
设$p$为电子的四维动量，$k$为光子的四维动量，$m$为电子质量，求下列矩阵的迹
\begin{itemize}
\item{$\slashed{k}$}
\item{$\slashed{p}\gamma^5$}
\item{$\slashed{p}\slashed{k}\frac{1}{\slashed{p}+\slashed{k}-m}$}
\item{$\slashed{k}(\slashed{p}+\slashed{k}+m)^3\frac{1}{\slashed{p}+\slashed{k}-m}(\slashed{p}-\slashed{k}+m)^2\slashed{k}$}
\item{$(\slashed{p}+m)\gamma^\mu\frac{1}{\slashed{p}+\slashed{k} -m}\slashed{k}\gamma_\mu\slashed{k}$}
\end{itemize}
\ech
\end{frame}


\begin{frame}
\chtitle{第5题(0.5分）}
\bch
电子和$\mu$子可以分别看成独立的两个Dirac场$\psi_e$，$\psi_\mu$的粒子。和电磁场一起的拉氏密度为
$$\lagr = -\frac{1}{4}\Fup\Fdown + \bar{\psi}_e(\ii \slashed{\partial} - m)\psi_e + \bar{\psi}_\mu(\ii \slashed{\partial} - m)\psi_\mu$$
画出散射过程
$$e^+e^- \rightarrow \mu^+\mu^-$$
的非零的最低阶近似Feynman图。
\ech
\end{frame}

\end{document}
